\chapter{Introduction \label{chapter:intro}}

This thesis presents an agent-based modelling framework as a step towards understanding the ``complex, integrated, tightly-coupled global fabric of exchange'' (\citealt{korowicz_tipping_2010} p.2) that makes up the modern spatial economy. The most pressing spatial economic questions require a way to explicitly model this `fabric of exchange'. For example, in adapting to climate change and rapidly shifting energy futures, how will changing costs affect economic welfare? Are the spatial components of these changes well understood? Once understood in more depth, can this knowledge be used to aid the development of a sustainable energy and climate future?

This can be framed as a broader theoretical question: how do changing space costs (the costs of moving people and goods across distance) affect the morphology of the spatial economic landscape and how does the resulting morphology feed back into those choices?

This thesis examines the theoretical factors of this problem in depth by comparing agent-based modelling and economic ideas, using the comparison to build its model framework and presenting results from that framework. Agent-based modelling is highly flexible but its theoretical foundations are still developing. The assumptions of more traditional economic theory are brittle but it has a deep well of tested, nuanced theoretical knowledge. By examining both in depth, the thesis hopes to contribute to modelling approaches that capitalise on the best of both worlds.


\section{Introductory discussion and aims}

\subsection{Agent-based modelling and spatial economics \label{sec:mot}}

In \gls{abm}, the `agents' are distinct code objects, programmed to interact with their environment and each other. The actor-centred nature of \gls{abm} makes it perfectly placed for investigating some of the most fundamental questions about spatial economics. It should, in theory, be ideal for tackling two of the most challenging problems that space imposes. 

First, it provides a way to model unique individual decisions and how they interact. The costs of moving people and goods across geographical distances change, as do the costs and benefits of proximity. As these changes take place, each agent makes many choices based on the unique set of factors their location imposes on them. Real-world agents have different locations; from where each of them stands, the costs they face differ. 
%A simple example: any three agents at different locations buying from each other must do so over different distances: no single market-clearing price is possible.

Secondly, these interactions are not bought together into neat auctions that happen at a set time and place. These two points together break market clearing assumptions central to many economic analyses (section \ref{sec:ace}). The aggregate result of all these decisions over time is written into the morphology of the landscape as a pattern of settlements and dynamic flows. Isard identified these vital qualities of the spatial economy in the 1950s. As he put it:
%Real-world agents have different locations; from where each of them stands, the costs they face differ.

\begin{quote}
``Each individual commodity market including its labour, capital and land orientation possesses a particular structure which offers a certain resistance to change. Some change frequently, others slowly. Some are active, others highly inactive. By definition those markets of a relatively permanent nature, of persistent inactivity, are grouped together as the essential ones... Their combined structure determines the basic structure of the gestalt whole, of the space economy under question.'' (\citealt{isard_location_1956} p.39)
\end{quote}

Using contemporary language, this `gestalt whole' is the emergent structure of the spatial economy - something \gls{abm} is, in theory, primed for. Yet while there is a large body of \gls{abm} work examining a range of spatial and economic ideas, as chapter \ref{chapter:ABM} discusses, \gls{abm} as a discipline seems to have stayed clear of tackling these more `general' spatial economic problems. 

To frame this as a question that might be suitable for an \gls{abm} research agenda: what happens to the spatial economy as a whole when costs change, given that space itself imposes heterogenous costs and rules out a single auction-style clearing market? This heterogeneity is \emph{the} essential element of space and a tempting target for agent modelling. \gls{abm} can enable agents to make spatial decisions outside of those assumptions, independently of others' actions.

This thesis presents a set of spatial economic agent models leading to a `transmission belt' model connecting supply and demand in a spatial setting. The goal is much narrower than Isard's `gestalt whole': it succeeds at producing a stable equilibrium of consumption and production, but the process of getting there involves a series of modelling simplifications and compromises. This process is a key focus for the thesis, used to ask: what obstacles have stopped \gls{abm} from tackling this classic spatial economic problem? The development of the thesis model is used to address this question in two ways.
%It does not attempt an `agent-based spatial theory of everything'. The goal is more mundane: to find out why \gls{abm} has mostly avoided a research area it seems so perfectly suited for. 

%Competition has proved the largest challenge. Some \gls{abm} theorists would argue that neoclassical economics does not solve it either, but assumes it away. I don't want to argue that: static assumptions can be highly useful.
Firstly, the goal for the thesis has been to start from first principles: to create a model that explores one possible route to a foundation for the broad question above. That requires starting `from the ground up': with agents able to make choices based on the costs they face in their particular time and place. The agents are kept as simple as possible to keep the focus on the more important factor (so it is argued) in this type of model: its detailed coding structure and the dependencies between components (see below). The separate components of supply and demand in a spatial setting are examined and methods used that best suit \gls{abm}. 

Secondly: while building up from these first principles, another central goal has been to recognise and \emph{use} the huge contribution made by older modelling methods and thinking, especially that of economics and the relatively new field of \gls{ge} that began with Paul Krugman's \gls{coremodel} in 1991 (section \ref{sec:coremodel}). \gls{ge} works in the `general' tradition of economics: it ``aims to explain the formation of various forms of economic agglomeration in geographical space using a general equilibrium framework'' (\citealt{fujita_frontiers_2005}, p.377). Though an array of spatial economic ideas have been used (see chapter \ref{chapter:location}), the thesis pays particular attention to \gls{ge}. It has two claims to be the most relevant modern approach to spatial economics, each of which makes \gls{ge} very useful for the thesis. First, a highly self-aware approach to model-building and, due to this, a deep sense of the history of its discipline. Second, a canonical model (the \gls{coremodel}) that attempts to boil space down to its absolute essentials, within the strictures of a neoclassical framework. The thesis `disassembles' the \gls{coremodel} to examine its workings and compare to the agent model built here (section \ref{sec:coremodel}). 
%The \gls{coremodel} of \gls{ge} is used as a comparator for understanding how the different modelling methods and philosophies of \gls{abm} and neoclassical economics impact on the way spatial models have been developed in those fields. This is done by `disassembling' the model components themselves as well as investigating the quite different ethos it has. 
%The goal is not to recreate \gls{ge} models in agent form. The aim is to look for the unique contribution that \gls{abm} may be able to make to spatial economics and what may have been holding it back from doing that. The agent model framework in the thesis has thus been created along a parallel track selecting concepts and ideas that worked, rejecting those that did not and testing new methods more suited to an agent structure.
%\gls{ge} embodies a unique combination of neoclassical requirements and high self-awareness, providing a very useful counterpoint for thinking through the issues in \gls{abm}.

Attempts to `convert' elements of \gls{ge} into an \gls{abm} approach do already exist; Christopher Fowler's work on this has been invaluable (see section \ref{sec:fowler}). But where Fowler has gone from the \gls{coremodel} towards \gls{abm}, the `direction of travel' for this thesis is different. The starting point is actors themselves: the costs they face from their point of view and the decisions they make given those costs. 
%Because of this, the emphasis throughout the thesis is on defining cost and decision functions at the `micro' level and working out how to build an interaction structure able to let actors work with those functions.

Using `first principles plus history' provides a way to carefully examine the components involved in building an agent-based spatial economic model - not just in terms of practical coding choices, though that is a vital part of it, but also in theorising what those components and the model as a whole are doing and what they are for. It is also a way to find a balance between the common \gls{abm} rejection of neoclassical economic ideas in favour of complexity (section \ref{sec:abmlit}; see below also) and what Mirowski describes as a ``respectful subordination to the neoclassical profession, discouraging anything that might have been perceived as wandering too far beyond the pale'' (\citealt{mirowski_kicking_2007}). 

Chapter \ref{chapter:ABM} examines relevant methodological and modelling philosophy issues using this historical approach - but the use of history has not been an abstract process. These issues have arisen in attempting to solve real modelling problems during the process of model development. Milton Friedman's 1953 paper on the methodology of economics (\citealt{friedman_essays_1953}) has been particularly useful. It has been dismissed by many social scientists and (as section \ref{sec:realism} discusses) some agent modellers, but they seem to have caricatured its meaning in the process. As Hoover notes, Friedman was interested in ``the quotidian practice of economics, not abstract epistemology'' (\citealt{maki_milton_2009} p.303) and this makes it an excellent resource. This thesis is written in the same spirit: it is about the `nuts and bolts' process of model-building. 

The rest of the introduction is organised as follows. Each of these sections provides cross-references to the rest of the thesis to give an overview of the upcoming material:

\begin{itemize}
	\item The following section summarises the above discussion into a set of aims and objectives, to be used as reference for the conclusion;
\item Section \ref{sec:themodels} explains the models presented in the thesis, discussing how they map onto the standard economic distinction between `partial' and `general' models. It also explains how the concept of `dependencies', borrowed from \gls{oop}, is used to frame the thesis argument;
\item Section \ref{sec:streetlight} sets the context for the comparison of \gls{abm} and other economic ideas by examining how different theories work under different `streetlights'; it also argues that \gls{abm} has a more acute problem than just the standard streetlight effect;
\item Section \ref{sec:problem} highlights some of the empirical reasons that more attention paid to agent approaches may benefit spatial economics, setting up a comparative discussion in the conclusion;
\item Section \ref{sec:org} gives a brief chapter summary and an overview of terminology.
\end{itemize}

\subsection{Aims and objectives \label{sec:aims}}

To summarise the above in the form of aims and objectives, this thesis' aims are as follows. Specific reference to these aims and objectives will be made in the conclusion (chapter \ref{chapter:conc}). The aims are:

\begin{itemize}
\item \textbf{Develop an agent-based spatial economic model framework in which agents with heterogenous locations can optimise welfare given changing costs}
\item \textbf{Use the development process of the model framework to answer the following: what obstacles are in the way of \gls{abm} exploiting its potential as a tool for spatial economic analysis?}
\end{itemize}

The following objectives are connected to the first aim:

\begin{itemize}
	\item Develop a spatial agent model framework and explain its workings in detail in the thesis;
	\item Compare agent-based and economic approaches to model-building and use this analysis to inform the thesis' model development.
\end{itemize}

The following objectives are connected to the second aim:

\begin{itemize}
	\item Identify obstacles revealed in modelling development at the coding level and a route through them;
	\item Identify obstacles revealed by the comparison of different approaches to model-building in \gls{abm} and economics, with specific emphasis on \gls{ge}.
\end{itemize}

The following chapters describe the route the model found through these obstacles at the coding level. The details of the obstacles themselves are raised throughout, and these are then bought together for a synthesised analysis in the conclusion. The second `obstacles' objective is more theoretical: the important differences in approaches to model-building, it will be argued, are about how \gls{abm} has developed a dysfunctional relationship with assumptions.
%The implications for model-building are far from abstract, however: they have been used to inform concrete model choices throughout. 

To pre-empt the conclusion, two key obstacles have been identified: a lack of research directly tackling the space of `dependencies' in \gls{abm} model development (see section \ref{sec:themodels} below); a broader sense that \gls{abm} has broken all ties with history and can learn nothing from past modelling efforts (see section \ref{sec:streetlight}). The conclusion analyses these obstacles in depth, building on what has been presented throughout the thesis, and presents an analysis of the model framework's dependencies and assumptions in order to highlight some possible ways forward.

\section{The models \label{sec:themodels}}

A series of agent-based spatial economic models are presented, built with a common framework. Chapter \ref{chapter:modelfw} explains this framework in detail. As well as a series of testing models (in the first sections of chapter \ref{chapter:models}) the thesis contains three main model set-ups. Note that capitalised `People' and `Firms' are used to describe model object representations, as per the \gls{oop} norm of capitalising objects, to distinguish them from `people' and `firms' generally. These three modelling scenarios are:

\begin{enumerate}
	\item `People' agents are mobile, able to optimise their locations as part of their larger economic decision set, with Firms' locations static (section \ref{sec:results}).
	\item `Firm' agents are mobile, able to seek revenue-maximising positions in relation to an immobile market of People (section \ref{sec:mobFirms}).
	\item A `transmission belt' model linking supply to demand (chapter \ref{chapter:connectedModels}): demand signals from People feed through to Firms. Using price signals, Firms are able to reach a market-clearing level of stock production that takes advantage of production-level economies of scale and allows for agents with heterogenous locations - thus keeping the `transmission belt' between supply and demand turning (see below for more on this).
\end{enumerate}

The first two model set-ups examining mobile People and mobile Firms are `partial' and the transmission belt model an attempt at a `general' analysis; the next section explains this distinction. 

%For the first two `partial' models, fixing the location of Firms and People allows the framework to be used to explore key issues and `dependencies' leading to the final model. 
What is the rationale for these two `partial' scenarios? They have been chosen because they provide the most insight for understanding assumptions and dependencies. During the process of model development, many options were tested. Some failed to work at all and some partially worked; a number of these are discussed in section \ref{sec:partialtogeneral}. The two presented in the thesis prior to the transmission belt model are the most informative examples.

The `mobile People' models achieve two things. Firstly, they are used for testing a `density cost' proxy. This is a very simple way to consider the effect of a range of proximity costs from congestion to land markets, very suited to an \gls{abm} approach as the cost is centred at each agent's own location. Section \ref{sec:results} explores its ability to produce a range spatial equilibrium outcomes, while section \ref{sec:extsinspace} provides some theoretical rationale. These spatial outcomes were a surprise, but a very useful one, providing an effective example for thinking about assumptions in \gls{abm}. 

Second, the `mobile People' models show that (within the model framework) commuting and the movement of goods cause the same agent reaction when changed. Equally, the `base' part of good costs (that is, without any spatial component) and `base' wage cause the same reaction as each other. `Commuting' is thus reframed as the spatial part of the wage, in the same way that goods have a spatial cost and a base cost. This idea is useful in the transmission belt model (see below).

In the `mobile Firms' model, two Firms attempt to maximise their revenue by optimising their location on a line of People. Creating Firms able to optimise good prices and wages was one of the problem areas for the framework; these models use the more transparent case of location optimisation to explore why this optimisation was so difficult. It uses an idea adapted from Jane Jacobs (\citeyear{jacobs_nature_2001}), framing Firms' decision process in terms of `\gls{dmr}' (see section \ref{sec:dmr}). The main lesson is that independently acting agents in a spatial setting can cause signalling confusion that make emergent spatial outcomes less likely. Using a canonical \gls{ge} utility function and a noise function, an analysis of success and failure is given. 
%is used to investigate the nature of information processing in a simple two-Firm scenario,

The transmission belt model presents one solution that allows agents to coordinate economic activity to a successful dynamic supply/demand equilibrium, capable of moving between production `regimes' where more-efficient producers can emerge. Section \ref{sec:prodfuncts} gives the production function used; section \ref{sec:abmprod} explains how problems of strategy were avoided in the coding of production. Reframing person-moving and good-moving costs into separate spatial and non-spatial elements helped by providing a rationale for People to buy goods directly with time (or labour) from a range of Firms. Thus a single distance cost, rather than two (distance moved by goods and commuting) can be varied to investigate its impact. The \gls{dmr} investigation leads to a novel way of allowing price signals to work that avoided a range of the thornier timing issues (discussed in chapter \ref{chapter:modelfw}). Leijonhufvud's concept of economic `laws of motion' (section \ref{sec:ace}) has been adapted, with the addition of oscillation damping.
%For the final transmission belt model, however, neither type of agent have direct mobility. People are able to shift their buying/working patterns from a static location, thus providing the model with the ability to change its morphology, and this can be seen as a movement of labour.

%These three sections are the main trunk of modelling work. Many branches presented themselves along the way, each involving other possible modelling choices that were not pursued in this thesis. (A number of these `paths not taken' are described throughout the results chapters.) This is one of central findings from the thesis work: within this way of thinking about \gls{abm}, there is clearly a very rich seam of further work.
\subsection{`Partial' versus `general' models}

General equilibrium models (of which the \gls{coremodel} is a geographical example) must link supply and demand to produce an equilibrium outcome that can be claimed to represent a whole economy. `Partial' models concentrate on sub-components of this, looking only at supply or demand and keeping much of the model fixed by assumption. 

Overman explains the importance of the difference between partial and general:

\begin{quote}
``Time and again economists have found that this is vitally important because partial equilibrium reasoning (just looking at a small part of any problem) often fails to provide the right reasoning in a general equilibrium context (when taking the economy as a whole)'' (\citealt{overman_can_2004} p.506).
\end{quote}

Birkin and Wilson define the difference between `partial' and `general'  \emph{location} models similarly:

\begin{quote}
``Partial theories are concerned with the location of a small number of production centres in relation to a given distribution of other centres and of markets. General location theory, on the other hand, aims to determine simultaneously an optimum distribution of all production centres and markets.'' (\citealt{birkin_industrial_1986} p.178)
\end{quote}

Combining `general' theories with the heterogeneity imposed by space is challenging. This problem has been usefully avoided in neoclassical economics through \emph{a priori} equilibrium assumptions. Analytic approaches can use models that assume an equilibrium endpoint has been reached. Agent models, however, are stuck trying to find a path to stable outcomes through interaction, unable to rely on those assumptions. \gls{abm} as a field has taken this coordination problem to be one of its main research aims: showing how `emergent' economic phenomena can arise through interaction. Fowler captures the tangle of relations involved in these economic interactions:

\begin{quote}
``the amount a firm can offer in wages is dependent on the amount it can produce and the price it receives for its goods. These quantities are affected, in turn, by consumer's wages, which depend on which firm employs them (and whether or not they have found employment at all.)'' (\citealt{fowler_taking_2007} p.277)
\end{quote}

The transmission belt model finds an agent-based way to link economic production and welfare in a spatial setting. Agents coordinate through decentralised actions. Production takes place and this feeds through into collective welfare. For the thesis, `welfare' refers to the aggregate level of utility agents are able to get. The term is useful for distinguishing discussion of model-wide optima in comparison to the utility of individual agents. The transmission belt model equilibriates supply and demand and can move between production equilibria as distance costs and utility are varied. In doing this, a range of varied prices can emerge based on differences in both location and production scale: no `market clearing price' is necessary for the equilibrium to sustain. As with the other models, however, there were compromises to achieve this end; these are discussed in-depth in the chapter presenting the model (chapter \ref{chapter:connectedModels}).
%This is as close as the thesis gets to a `general' model, in contrast to `partial' models where supply and demand are treated separately. Partial models are very useful: the first two main model results in this thesis are partial, and all applied models must to an extent be partial (an issue discussed in depth in section \ref{sec:appliedvstheory}).

\subsection{Dependencies \label{sec:deps}}

The thesis argues that model `dependencies' are unavoidable. In the process of attempting to move from a `partial' to `general' supply-demand link, the issue of dependencies arose naturally. Choices made in one area of a model determine those required in others. These choices are often mutually exclusive: one set can rule out another. This language of `dependencies' is taken from \gls{oop}. In \gls{oop} terms, dependencies are ``the villain of the piece''\footnote{Scitovsky (\citeyear{scitovsky_two_1954} p.144), discussing the role of externalities in neoclassical economics; see section \ref{sec:extsinspace}.}: the goal is to achieve `loose coupling' (see e.g. \citealt{jana_java_2005}) where objects are, as far as feasible, not reliant on any particular code choice elsewhere in a program. This philosophy, while useful for developing code itself, presents problems when applied to the process of modelling social systems.  As section \ref{sec:virtualworlds} discusses, it helps tilt \gls{abm} towards a `virtual world' philosophy (see below also). There is a sense that if good \gls{oop} practice has been followed, loosely coupled agents can interact as independent beings in their provided environment and a pseudo-empirical world \emph{in silico} will naturally result, leaving the modeller an observer-god of their own creation.

While an \gls{oop} concept, the language of dependencies is a very useful way to see the connection between older and more modern modelling methods. The spatial issues described above are, in reality, dependencies: spatial economics imposes a series of strictures on what modelling assumptions can be used (these are explained in section \ref{sec:spacesituate}). \gls{abm}'s error has been in thinking it can transcend these dependencies.

%They are not usually paid any attention in \gls{abm}. 
In economic approaches like \gls{ge}, these kinds of compromises, simplifications and choices are accepted and openly discussed. Section \ref{sec:coremodel} explains the \gls{coremodel} in detail; as with all economic models, it is a tightly knitted set of arguments built on mathematical links in a chain - all fully `dependent' on each other. It could be argued this is anathema to an \gls{abm} way of working. However, this thesis has provided an opportunity to explicitly examine how dependencies arise in the logical journey of model framework development. It has involved threading a path through these dependencies and understanding the modelling compromises required to do so. Rather than being an incidental obstacle on the path to creating an ideal \gls{abm} `virtual world', this thesis argues that these dependencies are an important part of the actual meaning of any model. 
%For example, the particular way that agent demand is coded requires an idiosyncratic structure on the supply side.
%While facets of the problem are solveable, each promising solution involves compromises in other areas. It is argued that this is especially important for \gls{abm} where `dependencies' are considered a sign that a model has not been properly implemented according to good \gls{oop} practice. 

Section \ref{sec:code_vs_function} makes a distinction between `descriptive' and `functional' mapping in models. Because of its deep roots in \gls{oop}, \gls{abm} has developed a implicit philosophy of mapping the real world descriptively onto code structures. As chapter \ref{chapter:ABM} argues, the pursuit of good \gls{oop} practice in framework development is one thing; applying the same \gls{oop} principles to the modelling philosophy itself quite another and may in fact damage \gls{abm} as a social modelling approach.
%The ability of agents to interact is indeed one of the key strengths of \gls{abm}. There is a balance to be struck, however, between leveraging the power of \gls{abm} and losing all analytic relevance.

By carefully breaking down the elements of an agent-based spatial economic model and comparing to the philosophy and methodology of \gls{ge}, this thesis hopes to show that the mapping of agent models is tangled into the minutiae of the code itself. \gls{oop} is still hugely powerful and, if carefully applied, can indeed be mapped onto the structures the modeller is interested in. But much of the mapping happens into the spaces in between the code where implicit meaning can exist.

These dependencies and compromises determined the form of each set of models. A clear examination of optimal Person choice versus optimal Firm choice requires fixing the locations of the other agent type in their respective models. By enabling and disabling certain aspects of the model framework, it is able to produce a series of spatial economic results that resemble classic spatial economic model outcomes (see section \ref{sec:litbits}) while also displaying novel behaviour.

%Key thing: again, this might take me a few runs to write down successfully. But: each of the model sections is used as an opportunity to highlight the dependencies between model choices. The final transmission belt models is the result of threading through a number of these dependencies, finding a set of model assumptions able to successfully allow the model to transition between different production equilibria.
%\subsection{`Partial' versus `general' models \label{sec:partialvsgeneral}}
%\subsection{The `transmission belt' between production and welfare}

This section ends with a note on the way the model framework has been used. The vast majority of the model outputs presented in the thesis run `live', with selectable parameters the user can interact with (see the appendix for more information on this). This is not incidental to the message of the thesis. Once a framework is tested enough to trust that it is robust, model behaviour can be quickly tested with this interactive approach and novel behaviour discovered. There is a cycle of experimentation and deduction. (This is discussed below also, in relation to the role of computation.)

There is a vital caveat to this. As section \ref{sec:generative} discusses, the strong \gls{abm} tendency to treat emergence and explanation interchangeably means that, often, agent analyses stop at the point a model appears to do something interesting. This should be just the start, however. The discovery of interesting model properties raises the question, \emph{why} is it happening? For the results presented here, this led to a deeper analyis of the underlying mechanism, often relying on a more analytic approach or creating simpler examples with complex assumptions stripped away.
%Not sure exactly how to frame this. But it ties to `generation vs explanation' and brushing under carpet of emergence. I never `brush under the carpet of emergence'. Each example displays emergent behaviour; finding out *why* it's happening is far from straightforward, and takes you back to first principles and toy examples to figure it out. But, there's a role for the manipulable model. I may have to drop that into the introduction.

To the extent that this found useful answers, these did not come solely from either the interactive agent framework, analytic calculations or more basic `toy' models, but an iterative combination of all of these, with each playing a key role. The role of visual interaction in \emph{discovering} potentially interesting dynamics was essential, however.

\section{The streetlight effect \label{sec:streetlight}}

\gls{abm} has positioned itself as the natural successor to what it sees as the `moribund Newtonian world-view' of neoclassical economics (see section \ref{sec:ace}). At present, however, \gls{abm} is stuck in a limbo of its own creation. Its methodology has played too strong a role in determining its research agenda. This is related to a common methodological problem called the `streetlight effect'. As Kirman explains, this ``corresponds to the behaviour of the person who, having dropped his keys in a dark place, chose to look for them under a street light since it was easier to see there'' (\citealt{kirman_whom_1992} p.134). The streetlight effect shapes what answers can be found as well as what questions can be asked. Maslow calls this `method-centring' in contrast to `problem-centring' (\citealt{maslow_psychology_1966} p.15); he summed it up in the now-famous aphorism: ``it is tempting, if the only tool you have is a hammer, to treat everything as if it were a nail'' (ibid). 

But \gls{abm} seems to have a particularly severe case of this problem - to a greater extent than the economic ideas it wants to supersede. Its methodology has not merely restricted what it can `see' but has heavily influenced its research questions. In combination with a prevalent rejection of everything neoclassical, this has left spatial economics surprisingly empty of agent-based analyses. 

In contrast, \gls{ge} is well-aware of its own limitations. \citeauthor{fujita_spatial_2001} explain that \gls{ge} works with - 
%There is a close link between the `strategic simplifications' essential to model and theory-building and the particular patches of reality that a model shines a light on. Criticising models for being simplistic makes little sense without first knowing what the model aims to achieve (when it may well be some element is indeed too simplistic, or other elements too complex). \gls{ge} theorists are particularly open about the strategic simplifications involved in the approach. 

\begin{quote}
``... assumptions that reflect not so much a realistic view of how the world works as a judgement about what will make the analysis of geographic issues manageable without doing too much damage to the relevance of that analysis.'' (p.6)
\end{quote}

More recently, Krugman - discussing his current approach to macro-economic modelling - notes that it:

\begin{quote}
``rests on some patently untrue assumptions about reality... models are an enormously important tool for clarifying your thought. You don't have to literally believe your model - in fact, you're a fool if you do - to believe that putting together a simplified but complete account of how things work... helps you gain a much more sophisticated understanding of the real situation. People who don't use models end up relying on slogans that are much more simplistic than the models.'' (\citealt{krugman_debt_2010})
\end{quote}

%Analytic economics and \gls{abm} are doing their searching under different streetlights. Many economists are well-aware of the limitations of their methods. 
Overman also acknowledges ``the tendency to privilege particular economic forces purely because they are more amenable to the theoretical and empirical tools used by mainstream economists'' (\citealt{overman_can_2004} p.504). Summer notes the result: ``it is all too easy to confuse what is tractable with what is right'' (\citealt{summers_scientific_1991} p.145). Krugman, as is often the case, puts it best: ``the methodology of economics creates blind spots. We just don't see what we can't formalise'' (\citealt{krugman_how_2005}). The result is that clear empirical realities often find themselves ignored if they fail to fit into existing theory. \gls{ge}'s openness about this is one of its main strengths: it imbues the approach with a sense of caution. 

%In contrast to the relative abundance of self-awareness in economics, the limitations of \gls{abm} are less often acknowledged.
\gls{abm} theorists, on the other hand, often argue they are immune from the streetlight effect by default. The promise of a much deeper realism aids this sense that it has broken free of the `mathematical straightjacket' (\citealt{martin_new_1999}) of neoclassical theory. However, this thesis argues that \gls{abm} has developed a unique straitjacket of its own making. The methodology of \gls{oop} itself, as well as the way \gls{cas} theory perfectly marries to \gls{oop}, are at least partly the cause. This is not a new thought: to mix metaphors, some argue the Santa Fe Institute has been overly dominant, ``retailing the snake oil of `chaoplexity' '' (\citealt{mirowski_kicking_2007}, p.361 paraphrasing \citealt{horgan_end_1998}). That is going too far: the ideas within the \gls{cas} canon have clearly opened up huge new research avenues and continue to be fruitful. The complexity agenda has been driven by a desire to understand social systems through treating them as composed of interacting parts. However, in \gls{ace} (an economic sub-discipline of \gls{abm}; see section \ref{sec:ace}) this has dampened the pursuit of many of the traditional economic research problems, leading to it carving out a quite separate theoretical territory. 

%Whether this mix of \gls{oop} and \gls{cas} is solely responsible for the `year zero' flavour to \gls{abm} is hard to say.
Two recent examples highlight this notion that \gls{abm} has the power to declare a theoretical `year zero'. Reflecting on the nature of agent models, one concludes that \gls{abm} is coming into its own in a time when ``the classical canons of scientific inquiry melt away into the vestiges of history'' (\citealt{batty_generic_2012} p.4). Another argues:

\begin{quote}
``half a century ago, the idea of a model was in its infancy. Scientific theory essentially was based on formal and systematic theories, often represented mathematically, whose testing was confined either to controlled experiments in the laboratory or to various categories of thought experiment. Computation changed all that. The idea that a scientific theory could then be translated into an intermediate form - called a model - represented a way of enabling controlled experiments to be carried out not on the actual system of interest but on a computable abstraction of that system.'' (\citealt{batty_perspectives_2012} p.21)
\end{quote}

The idea of this kind of model goes back to at least the turn of the twentieth century, and possibly to Bacon\footnote{The Oxford English Dictionary has references from 1901 for `model' as ``a simplified or idealized description or conception of a particular system, situation, or process, often in mathematical terms, that is put forward as a basis for theoretical or empirical understanding, or for calculations, predictions, etc.; a conceptual or mental representation of something'' (\citealt{oed_oxford_2002}) }. But this notion that a \emph{computational} model is something historically unique ties to the belief that computers have allowed modelling to transcend that history. It is irrefutably true that computational modelling has radically transformed research. As Berlinski puts it:

\begin{quote}
``The computer has changed the very nature of mathematical experience, suggesting for the first time that mathematics, like physics, may yet become an empirical discipline, a place where things are discovered because they are seen.'' (Berlinski 1997, quoted in \citealt{gold_proof_2008} p.36).
\end{quote}

%For the work presented in this thesis, computation and agent modelling offered a way to create a virtual experiment station. By itself, this did not reveal new insights, but it did show where to look for them. 
Beyond the hyperbole of this quote there is a useful seed of truth. This is precisely the point made above about the cycle of experimentation and deduction: computation has made that possible. For \gls{abm}, though, this power can be a trap, precisely because it empowers the modeller to create very plausible-seeming `virtual worlds' (section \ref{sec:virtualworlds}) that are both compelling for the researcher and appealing to those outside of the discipline. In its place, compelling or appealing \gls{abm} work are not bad things - but they must not be confused with the scientific job that models can do. 

As section \ref{sec:virtualworlds} discusses, Friedman saw exactly the same tendency in his day: a clash between attempts at `photographic reproduction' versus `engines of analysis' (\citealt{friedman_essays_1953} p.35). This thesis argues that the fundamental nature of models has not changed since Friedman wrote this. The classical canons of scientific inquiry are most emphatically not melting away. This is a positive thing for \gls{abm}: as this thesis explores, it means that the rich well of thinking about modelling in economics is still highly relevant. 
%At root, there is a confusion between models on the one hand and modellers' understanding on the other. The latter should be informed by the former; the relationship between them is not simple, however. 

The models in this thesis are only asked to pass `face validation' (\citealt{oeffner_agent-based_2008} pp.35): does it work as a self-consistent set of relations, able to explore the question of how changing costs affect actors in a heterogenous, abstract geography? Is it a useful tool for ``probing the internal consistency of a theoretical position?'' (\citealt{di_paolo_simulation_2000} p.1). The impact of these differences in modelling philosophy affect every facet of model building, however. For this thesis, they impact most keenly on the way dependencies are considered and how they relate to assumptions about reality. The issue is examined by comparing how agent modellers and earlier economic theorists have thought about what they are doing. What do the two approaches think their models are `explaining'? How do they think about simplicity and realism? The way model assumptions are theorised is key. There are very few theorists without a personal pet-hate assumption that they see as - in Hayek's phrase - self-evidently disregarding ``everything that is important and significant in the real world'' (\citealt{hayek_use_1945} p.530). No assumption can be dismissed so easily since what is `important and significant' is in the eye of the beholder. In assessing assumptions, ``everything depends on the problem'' (\citealt{friedman_essays_1953} p.7).

%The motivation for doing this does not come from a desire to add new knowledge to the philosophy of modelling. The goal is instead to offer the suggestion that there are fertile areas of agent modelling currently unexplored, and to build a simple model framework to illustrate possible ideas. 
%So: the point of this bit of `model' writing? 
\section{The spatial economic `problem' \label{sec:problem}}

%\gls{ge} has its origins in an empirical misfit between older trade theory built on factor difference and the fact that international trade often consists of the exchange of similar goods (section \ref{sec:cmstory}). Neither model approach is necessarily `wrong'; the point is finding something that helps shed light on an aspect of reality previously beyond the reach of analysis. This section outlines some of the empirical motivation for attempting to assess the value of \gls{abm} for spatial economic analysis.

As discussed above, many theorists of have argued for the pursuit of a `general theory' of the spatial economy capable of understanding the `gestalt whole' of its interactions. While the focus of this thesis is methodological, not empirical, this section outlines some of the empirical motivation for pursuing new spatial economic ideas. 

Spatial economics has developed during a particular historical period of constantly dropping costs through both an increase in production efficiency and a drop in transport costs. Modelling approaches have been shaped by this, sharing much the same general view of space costs as Glaeser and Kolhase, who have claimed:

\begin{quote}
``reduced costs, and the declining importance of the good-producing sector of the economy, means that in our view it is better to assume that moving goods is essentially costless than to assume that moving goods is an important component of the production process'' (\citealt{glaeser_cities_2004}, p.199).
\end{quote}

They also argue that ``there is little reason to doubt that this decline will continue'' (\citealt{glaeser_cities_2004}, p.197). This assumption of space costs inevitably dropping is a widespread part of the literature. Lang points out that \gls{ge} models aim to:

\begin{quote}
``examine how the spatial distribution of economic activity changes as transportation costs slowly decline. The slow decline is meant to mirror the actual decline in transportation costs observed over the modern era of civilisation.'' (\citealt{lang_heterogeneous_2010} p.191)
\end{quote}

%note: orig paper for 'history' quote is krugman, venables 2005, globalisation and the inequality of nations
As Neary notes (\citealt{neary_hype_2001} p.536), this is a bold claim of telling the `history of the world part one', spanning ``first caravels, then steamships and railroads, then air freight...'' (\citealt{fujita_spatial_2001} p.253) This is a somewhat `Whig history' approach that assumes the direction of costs is inevitably one of decline -  but is that a problem? 

Twenty-first century challenges are breathing new life into spatial economic questions. Theorists are rediscovering that spatial costs may go up as well as down. During the next few decades an energy revolution must take place to avoid the worst effects of climate change. What price must carbon be to keep within a given global temperature? How long will any switch to new infrastructure take? (\citealt{kramer_no_2009}; \citealt{jefferson_accelerating_2008}) Will the cost of fossil fuels impact on the spatial economy as much as climate change? (\citealt{wilkinson_peak_2008}; \citealt{bridge_geographies_2010}; see \citealt{sorrell_assessment_2009} for a thorough analysis of the issue of oil depletion). More broadly, what role might a deeper understanding of spatial economics play in reshaping a post-carbon economy?

The more qualitative approach of `economic geography' examines many of these questions. The idea of `transition' - the need to consciously move to a low-carbon economy through planning and civic engagement (\citealt{hopkins_transition_2008}) - has gained a great deal of research territory in recent years. But it brings many assumptions with it. Economic analysis is eschewed; in fact it is common to find a polar view of an `` `Industrialised World' of standardised-generic production... associated with commercialism, efficiency and branding'' versus an `` `Interpersonal World' of specialised-dedicated production... associated with trust, local renown, and spatial embeddedness'' (\citealt{morgan_worlds_2006} p.22). Many geographers have already taken sides in a battle where localities fight to `respatialise' social, economic and environmental goods lost to globalisation (\citealt{glasmeier_book_2007}). 

While many see `relocalisation' as the self-evidently correct normative response to a perceived future of rising space costs, the most common view from the economics profession is summed up nicely by Harford: `buying local is simply a lifestyle choice - not the difference between environmental salvation or damnation' (\citealt{harford_frequent_2007}). The problem with accepting Harford's argument at face value is - as discussed in detail in this thesis, and argued by Paul Krugman (see section \ref{sec:spacesituate}) - economists have been spectacularly good at ignoring geography, because the `dependencies' involved in modelling space makes many central economic assumptions unworkable. As section \ref{sec:spacesituate} points out, critiques of this ``wonderland of no dimension'' (\citealt{isard_location_1956} p.26) go back at least to Isard.
%ABM should be positioned to represent all the many agents and how they interact in space. But the problem is, in fact - as the history of spatial economics demonstrates - not easy.

`Transition' analyses argue that rising oil prices will have severe impacts, both spatial and economic. Economic takes on oil price shocks, though in-keeping with the economic avoidance of space, raise questions about this assumption. Oil price \emph{volality} may be more important than price change itself (\citealt{bachmeier_should_2008} p.528; see section \ref{sec:spacerisk} also), though, matching an argument often also used by transition thinkers, recent oil price effects on economic output may by ``caused by strong demand confronting stagnating world production'' (\citealt{hamilton_causes_2009} p.215). The oil shock economic literature also offers alternative theories to `transition' accepted truths. For example, in the transition literature, it is commonly suggested that oil price rises correlate to recessions, thus indicating its central role as the lifeblood of the economy. However, Segal suggests that the oil-price/recession connection has been driven mostly by monetary policy response: raising interest rates when high oil prices feed through to inflation, thus helping to cause recession without necessarily addressing the root cause of the problem (\citealt{segal_oil_2011} p.169-170).

%Can `transition' type questions have a quantititative answer?
A quantitative spatial modelling approach may help shed some light on these different views - and indeed, there are relevant quantitative examinations of aspects of the `transition' problem. For example, Kerschner and Hubacek, in a very rare occurrence, actually take on the `oil depletion' question directly (\citealt{kerschner_assessing_2009}). The MARKAL model approach (see e.g. \citealt{seebregts_energy/environmental_2001}) attempts to integrate national energy economics into one framework. As with the oil shock literature, however, these avoid any of the fundamental distance effects that changing space costs cause. This throws into sharp relief the fundamentals of spatial economics being missed: it is a web of interconnected agents making decisions separated by both time and space, in which changing costs can alter morphology, feeding back on those costs to create Isard's `gestalt whole'. \gls{ge}, aware of its own limitations, makes no claim in this regard. No \gls{abm} analysis directly addresses it, leaving these questions to qualitative geographers. This thesis hopes to open up some space for exploring quantitative approaches to the issue using agent modelling. This will be returned to in the final section examining future work (section \ref{sec:theend}) where it will be asked: is \gls{abm} well-positioned to deal with understanding how the `gestalt whole' spatial economy changes?
%Whether or not the polarity of cost changes are important for the shape spatial models have taken, these


%So: some theorists think transport costs are almost irrelevant and will become more-so. Others simply do not deal with the geographical part of economic activity. Others take it as given that ``our basic local needs have become dependent upon a hyper-complex, integrated, tightly-coupled global fabric of exchange'' (\citealt{korowicz_tipping_2010} p.2). They tend to argue that resilience should be the goal and assume that localisation means more resilience - but offer no routes into testing those assumptions through modelling the problem.
%What the thesis model tries to do is simple: ask how changing space costs affects both morphology and welfare, given that the feasible distance of moving goods and people may change. It does this by using a set of economic relations between actors to  investigate how they react to cost changes. It works with the fact that space itself imposes heterogenous costs: if actors have different locations, the costs they face must differ, and the outcomes of their decisions may leave costs differentiated.


\section{Organisation of the thesis \label{sec:org}}

\subsection{Terminology}

A glossary and list of acronyms is provided after the table of contents. The first time the glossary terms appear in the thesis, they will be in bold. The first time acronyms are used, the full phrase will precede them.

A brief note on the terminology used to describe fields of research. `Spatial economics' is used as a top-level heading to describe all theories that attempt to say something about how economic activity and geography interact. \gls{ge} has already been mentioned: it is a specific approach beginning with Krugman. As Brakman \emph{et al.} say, it is ``an attempt to put more geography into economics'' (\citealt{brakman_introduction_2011} p.xxii). \gls{ge}'s other common name, the `new economic geography', is no longer really apt since it is twenty years old now (ibid). Confusingly, `geographical economics' is completely different to `economic geography': the latter a much broader church, often qualitative, case-study-based and empirical, and sometimes openly hostile to \gls{ge}. Almost all current academic work on `transition' issues comes under the `economic geography' banner. This point is discussed more in section \ref{sec:spacesituate}.

\subsection{Chapter outline}

Chapter \ref{chapter:location} explains the key economic and spatial ideas used in the thesis. It presents the utility and production functions chosen to build the model. The theoretical context for the economic and spatial ideas used is given. The \gls{coremodel} is given its own section. The way in which the thesis theorises the costs of moving people and goods\footnote{Glaeser talks of cities eliminating ``transport costs for goods, people and ideas'' (\citealt{glaeser_cities_2008} p6.). This thesis sticks to goods and people.} across space is presented. The chapter ends with discusion of some strategic simplifications.

Chapter \ref{chapter:ABM} has two aims. Firstly, it introduces the idea of \gls{abm} in more detail. Secondly, it looks critically at how \gls{abm} has developed and puts it in the context of some key ideas from more traditional economic modelling theory. The chapter ends by discussing the use of utility and production in an agent modelling context.

Chapter \ref{chapter:modelfw} explains how the model framework has been put together. This covers: how the model setup works; how actors have been set up; an explanation of the model scenarios used; finishing with a more detailed discussion of some of the key framework areas that need more background detail. 

Chapter \ref{chapter:models} firstly presents a series of testing models to demonstrate the framework achieves its fundamental economic outcomes. It then presents the two `partial' model set-ups, examining `mobile People' and `mobile Firms'.

Chapter \ref{chapter:connectedModels} presents the transmission belt model. It begins by discussing how the model resulted from a series of strategic choices; this section also outlines some of the approaches that are excluded by the the particular set of dependencies the model has. The transmission belt model is then put through a series of runs, each introducing successive elements: a `damping' method for successfully finding prices; non-spatial models showing the equilibrium outcome; an examination of the interaction of `love of variety' and `economies of scale'; and a series of models introducing a spatial element.
%Third, a series of `monocentric' models are used to examine how space cost change affects model outcome. Fourth, a series of models are presented that examine economic coordination between actors. The chapter uses a mix of model outputs and analytic results to explore the findings.

Chapter \ref{chapter:conc} concludes by discussing the aims and objectives in the light of the thesis model results. It then examines the set of specific dependencies and assumptions revealed through the model development process and uses them as the focus for highlighting both the limitations of \gls{abm} and new opportunities these limitations reveal.

The appendix provides a detailed overview of the model code structure and a full description of the downloadable model location and contents. 

Next, chapter \ref{chapter:location} outlines the economic and spatial ideas used in building the model. 

